\documentclass[11pt,a4paper]{article}
\usepackage[T1]{fontenc}
\usepackage{lmodern,amsmath,amssymb,amsthm,mathtools,booktabs,microtype}
\usepackage[margin=26mm]{geometry}
\usepackage[colorlinks=true,linkcolor=blue,citecolor=blue,urlcolor=blue]{hyperref}
\setlength{\emergencystretch}{3em}
\newtheorem{theorem}{Theorem}[section]
\newtheorem{lemma}[theorem]{Lemma}
\newtheorem{corollary}[theorem]{Corollary}
\theoremstyle{remark}\newtheorem{remark}[theorem]{Remark}
\newcommand{\Q}{\mathbb Q}\newcommand{\Z}{\mathbb Z}
\title{Recovery of elliptic curve equations from complete ratio equations}
\author{}\date{15 November 2026}
\begin{document}\maketitle
\begin{abstract}
Complete elliptic ratio equations, together with the secant and tangent
relations of a reciprocal point pair, imply both point-curve equations
provided that the two reference ratios do not vanish simultaneously.
Translation by a fixed point outside the two-torsion subgroup supplies
this nonvanishing without a rank hypothesis. For the quartic construction
under ordinary rational $abc$, the implication removes two rational
quadratic equations while preserving every solution, including identity
and torsion cases. The resulting 78-variable polynomial is a sum of
68 primitive quadratic squares with 1113 monomials. Its degree, coefficient
height, rational input projection and defining parameter threshold are
unchanged. No decrease in the number of variables is asserted.
\end{abstract}

\section{A polynomial implication for a reciprocal point pair}
Throughout this section let $k$ be a field of characteristic zero and
$D\in k^\times$. Use coordinates
\[
 f,h,g,j,s,F,G,L,M,S,T\in k.
\]
The letter $j$ denotes a coordinate, not an elliptic invariant.
Consider the equations
\begin{align}
 F=f^2,&\qquad G=g^2,\label{eq:squares}\\
 s(f-g)&=h+j,\label{eq:secant}\\
 s(h-j)&=1-D(F+fg+G),\label{eq:tangent}\\
 hL+fS&=0, & (1-DF)L+hS&=0,\label{eq:ratio1}\\
 jM+gT&=0, & (1-DG)M+jT&=0.\label{eq:ratio2}
\end{align}
These are quadratic equations; zero coordinates are allowed.

\begin{theorem}\label{thm:implication}
If \eqref{eq:squares}--\eqref{eq:ratio2} hold and $(S,T)\ne(0,0)$,
then
\[
                  h^2=f-Df^3,\qquad j^2=g-Dg^3.
\]
Consequently the additional equation $h^2-f+DfF=0$ is redundant on
this entire locus.
\end{theorem}
\begin{proof}
Put
\[
 C_U=h^2-f+Df^3,\qquad C_V=j^2-g+Dg^3,
\]
and denote the left sides of the secant and tangent equations by
\[
 r_1=s(f-g)-h-j,\qquad r_2=s(h-j)-1+D(F+fg+G).
\]
The following are identities in the polynomial ring on all displayed
coordinates:
\begin{align}
 h\{(1-DF)L+hS\}-(1-DF)(hL+fS)
 &=SC_U+DfS(F-f^2),\label{eq:id1}\\
 j\{(1-DG)M+jT\}-(1-DG)(jM+gT)
 &=TC_V+DgT(G-g^2),\label{eq:id2}\\
 C_V-C_U+(f-g)r_2-(h-j)r_1
 &=D(f-g)\{F-f^2+G-g^2\}.\label{eq:id3}
\end{align}
The equations therefore give $SC_U=0$, $TC_V=0$ and $C_U=C_V$.
If $S\ne0$, the first equality gives $C_U=C_V=0$. Otherwise $T\ne0$,
and the second gives the same conclusion. Only the specified nonzero
reference value is divided by; none of $f,g,h,j,L,M$ is assumed nonzero.
Using $F=f^2$ proves the last assertion.
\end{proof}

The proof is an identity argument over a field, rather than a statement
restricted to a dense open set. In particular it includes $f=h=0$,
$g=j=0$, tangencies, finite torsion points and vanishing ratio values
whenever they satisfy the retained equations.

\section{Uniform rescaling of the seven point coordinates}
The square coordinates need not be rescaled by the squares of the
factors used for the point coordinates. Rescaling all seven coordinates
by the same unit gives a quadratic presentation that is useful when
reference functions have variable denominators.

\begin{theorem}\label{thm:scaling}
Let $v\in k^\times$ be a further coordinate whose nonvanishing is
guaranteed by retained equations. In the system of Section~1 replace
the old coordinates by
\[
 (f_*,h_*,g_*,j_*,s_*,F_*,G_*)
       =v(f,h,g,j,s,F,G),
\]
and leave $L,M,S,T$ unchanged. The following quadratic equations are
equivalent to the original equations under this invertible substitution:
\begin{align}
 vF_*&=f_*^2,&vG_*&=g_*^2,\label{eq:scaled-squares}\\
 s_*(f_*-g_*)&=v(h_*+j_*),\label{eq:scaled-secant}\\
 s_*(h_*-j_*)&=v^2-D(vF_*+f_*g_*+vG_*),\label{eq:scaled-tangent}\\
 h_*L+f_*S&=0,&(v-DF_*)L+h_*S&=0,\label{eq:scaled-ratio1}\\
 j_*M+g_*T&=0,&(v-DG_*)M+j_*T&=0.\label{eq:scaled-ratio2}
\end{align}
If a point-curve equation is retained, its equivalent quadratic form is
\begin{equation}\label{eq:scaled-curve}
                         h_*^2=vf_*-Df_*F_*.
\end{equation}
On $(S,T)\ne(0,0)$ this last equation is a consequence of
\eqref{eq:scaled-squares}--\eqref{eq:scaled-ratio2}.
\end{theorem}
\begin{proof}
Substitute $f=f_*/v$, and similarly for the other six coordinates,
in each original row. Multiply each square-definition row, the curve
row and the two point-pair rows by $v^2$. Multiply each complete
ratio row by $v$. The resulting equations are precisely those
displayed. Their monomials have total degree at most two in all
retained coordinates, including $v$ and the reference ratios.
All multipliers are nonzero. Division by $v$ gives the inverse map
on every solution, and Theorem~\ref{thm:implication} proves the last
assertion by transport through this bijection.
\end{proof}

In particular, the new stored square is
$F_*=vf^2=f_*^2/v$, rather than $f_*^2$.
This distinction is necessary for the displayed degree assertions.

Suppose in addition that $k/\Q$ is quadratic and that $v^{-1}$ is a
retained linear function, as in the three-coordinate reference chart
of~\cite{Three}. The old difference abscissa is
\[
                    x(U-V)=s^2+D(f+g).
\]
It may be recovered while keeping its trace equation quadratic by
retaining one further field coordinate $r=s$ and imposing
\[
 s_*=vr,\qquad
 \tau=\operatorname{Tr}_{k/\Q}
                 \bigl(r^2+Dv^{-1}(f_*+g_*)\bigr).
\]
Both equations have degree at most two. A field coordinate here means
two rational coordinates. This is an exact recovery construction;
it does not remove those two rational coordinates or assert a reduction
in any complete defining polynomial.

\section{The reference nonvanishing}
Let
\[
                    E:y^2=x^3-Dx
\]
over $k$, and fix $P\in E(k)$ with $2P\ne O$. Suppose that $H,H-P$
are finite and have nonzero abscissas. Define
\[
             S=\frac{y(H)}{x(H)},\qquad
             T=\frac{y(H-P)}{x(H-P)}.
\]

\begin{lemma}\label{lem:nonvanishing}
For every such $H$, one has $(S,T)\ne(0,0)$.
\end{lemma}
\begin{proof}
If both values were zero, both ordinates would be zero. On the given
Weierstrass equation each finite point with ordinate zero is fixed by
negation and therefore belongs to $E[2]$. Hence
$P=H-(H-P)$ would belong to the subgroup $E[2]$, contrary to $2P\ne O$.
\end{proof}

For the construction below, $k=\Q(\rho)$, $\rho^2=2$,
$D=11-6\rho$, and
\[
 P=(a,b)=\left(\frac{57}{8}-\frac72\rho,
                  -\frac{165}{8}+\frac{417}{32}\rho\right).
\]
Here $b\ne0$, so $P\notin E[2]$ without any group-rank hypothesis.
The complete five-coordinate reference chart of~\cite{Reference}
recovers precisely such $H,H-P$ on every field solution. Its exceptional
leading-coefficient cases are included in that theorem. Thus
Lemma~\ref{lem:nonvanishing} applies to every solution of the retained
reference equations.

\section{The 78-variable quartic with 68 squares}
We use the complete quadratic family of~\cite{ABC78}, including its
counting convention that the input coordinate $X$ is included in the
total. Its rational coordinate list is kept unchanged. The coefficient
parameter $B$ is fixed and is not quantified.

Let $q_{B,1},\ldots,q_{B,70}$ denote its primitive integral quadratic
rows in their published order. Let $J$ consist of all row indices except
the two rational components of
\begin{equation}\label{eq:deleted}
                         h^2-f+DfF_U=0.
\end{equation}
The two rows have names \texttt{pair.cubic.0} and
\texttt{pair.cubic.1} in the exact source table. Define
\begin{equation}\label{eq:polynomial}
                       Q'_B=\sum_{i\in J}q_{B,i}^{\,2}.
\end{equation}
This gives a finite integer coefficient specification, with no
normalization or modification of any retained row.

\begin{theorem}\label{thm:quartic}
For every fixed coefficient parameter $B$, the old and new quadratic
systems have exactly the same rational zero set. For every positive
integer $B$, the polynomial \eqref{eq:polynomial} has 78 variables,
68 quadratic summands, total degree exactly four, integer content one,
and exactly 1113 coordinate monomials. Its coefficient height is
\[
 H(Q'_B)=\max\{C_*,B^2+C_0\},
\]
where
\begin{align*}
 C_*&=116167112798432587389449032604886634366655869137034483200,\\
 C_0&=1747145497790223592810060466409331738555986526175317640.
\end{align*}
Under the exact arithmetic inputs of~\cite{ABC78}, including ordinary
rational $abc$, every sufficiently large fixed odd $B$ gives
\[
 x\in\Z\quad\Longleftrightarrow\quad
       \exists\boldsymbol y\in\Q^{77}:Q'_B(x,\boldsymbol y)=0.
\]
The same sufficient threshold as for the parent family works.
No effective numerical threshold is supplied.
\end{theorem}
\begin{proof}
Every old solution is a new solution because only equations have been
removed. Conversely, the reference chart on a new solution still
recovers $H,H-P$ as above, so $(S,T_H)\ne(0,0)$ by
Lemma~\ref{lem:nonvanishing}. Apply Theorem~\ref{thm:implication} with
\[
 (j,F,G,L,M,S,T)=(k,F_U,F_V,L,L_0,S,T_H).
\]
All of its hypotheses are retained literal rows. It recovers
\eqref{eq:deleted}, hence both of its rational components. This proves
equality of the entire rational zero sets without changing any parameter.
Over $\Q$, a sum of squares is zero exactly when every square's base
is zero, so the two polynomials have the same zero set as well.

The variable list has 78 entries and exactly two of 70 rows were removed.
The only row containing $X$ remains $z_2-X^2$. Consequently the
coefficient of $X^4$ in \eqref{eq:polynomial} is one, independently of
$B$, and no other row contributes to it. This proves exact degree four
and primitive integer content.

The complete collected coefficient table specified in the next section
has 1113 coordinate monomials. Deleting \eqref{eq:deleted} changes only
parameter-independent coefficients. Every coefficient profile involving
$B$ is unchanged from the parent, as is the constant profile $B^2+C_0$.
The maximum absolute value of the remaining constant profiles is $C_*$.
The table and the explicit inequality checks below certify that these
profiles give precisely the asserted height for all $B\ge1$.
In particular this is an all-parameter coefficient statement, rather
than a computation at one test member.

Finally the parent's input-projection theorem applies because its full
zero set has just been recovered. Its arithmetic premises are retained:
the specified elliptic group and index statements, the fixed-data
simultaneous prime theorem and preparation conditions, the prepared
congruent-number simultaneous-twist realization, and ordinary rational
$abc$. No new prime-distribution or pole-realization assumption is used.
\end{proof}

\section{Exact specification and independent verification}
All coefficient data are in
\begin{center}\small
\path{computation/quartic/abc78_reduced_v40}.
\end{center}
The source table is authenticated by SHA-256
\begin{center}\small
\texttt{ceffdf6e0a38ac1c8c42812f5e7bef599231dc5bc4b5f7c4527303b81dfc11bb}.
\end{center}
The files \texttt{quadratic-family.json} and
\texttt{polynomial-family.json} contain every coordinate, row and
collected quartic coefficient. The separate checker does not import the
producer. It verifies identities \eqref{eq:id1}--\eqref{eq:id3},
authenticates all retained source rows, independently expands their
squares, and checks degree, support, primitive content and height.

For clarity, its all-parameter height check has the following form.
Set $m=\lfloor\sqrt{\max(0,C_*-C_0)}\rfloor+1$.
Every parameter-dependent coefficient other than the constant coefficient
is of the form $c_0+c_1B$. Its exact table satisfies
\[
 |c_1|(m-1)+|c_0|\le C_*,\qquad
 |c_1|m+|c_0|\le m^2+C_0,\qquad |c_1|\le2m.
\]
The first inequality controls $1\le B<m$; the latter two imply that
$B^2+C_0-|c_1|B-|c_0|$ is nonnegative and increasing for $B\ge m$.
Thus no parameter-dependent coefficient exceeds the stated height.
Both its attaining profiles occur in the table. Nonconstant profiles
have coefficients of one sign, so none vanishes for positive $B$;
this also certifies the ordinary monomial count for all such $B$.

The computation verifies finite identities and coefficients. It does not
constitute a new audit of the source's infinite arithmetic premises.
The full-set algebraic equivalence above is proved independently of
those premises; only the interpretation as a definition of $\Z$ uses them.

\begin{thebibliography}{9}
\bibitem{Reference}\emph{A complete five-coordinate elliptic reference chart},
project manuscript, corrected edition 14 November 2026,
\nolinkurl{papers/elliptic_charts/compressed_reference_20261114_v27.tex}.
\bibitem{ABC78}\emph{A quartic in 78 variables from a mixed-power height estimate},
project manuscript, 15 November 2026,
\nolinkurl{papers/quartic_existence/quartic_abc78_20261115_v34.tex}.
\bibitem{Pair}\emph{Trace heights on quadratic twists with $j=1728$},
project manuscript, Section 2,
\nolinkurl{papers/trace_heights/trace_on_all_points_20261113_v25.tex}.
\bibitem{Three}\emph{A three-coordinate quadratic chart for an elliptic
point and its translate}, project manuscript, 15 November 2026,
\nolinkurl{papers/elliptic_charts/three_coordinate_reference_20261115_v40.tex}.
\end{thebibliography}
\end{document}
